% Conclusion.

\section{Conclusion}
\label{sec-conclusion}

Nous avons appliqué au modèle de Heston la méthode de \citet{horvath2021} : apprendre une fois pour toutes la
fonction qui associe aux cinq paramètres la surface de volatilité, puis calibrer avec le réseau à la place du
pricer de Fourier. Elle fonctionne.

\paragraph{Le réseau.} Un petit réseau suffit. Le $2\times128$ est le plus précis des réseaux assez rapides : les
plus grands sont plus précis mais trop lents, et ajouter des couches n'aide pas. Le choix de l'entraînement compte
autant que la taille : avec Adam, l'initialisation de He et l'activation GELU, l'erreur passe de \num{5.09} à
\num{2.67}\pb\ sur le test, pour le même temps de réponse, plus de 70 fois plus court que celui de
Fourier en Julia. Ces choix ne s'additionnent pas : le meilleur réglage d'un hyperparamètre dépend des autres.

\paragraph{Réseau ou apprentissage classique.} Sur un petit jeu de surfaces complètes, un processus gaussien est
plus précis que le réseau. Le réseau l'emporte pour deux raisons : son erreur baisse encore quand on ajoute des
données, et la perte masquée lui permet d'apprendre sur toutes les surfaces, y compris celles qui ont un point
manquant. Avec toutes les données, il est 2,7 fois plus précis que le GP et son erreur maximale est cinq fois plus
petite.

\paragraph{La calibration.} Sur 100 surfaces simulées, le réseau retrouve les paramètres à moins de
\SI{0.22}{\percent} de leur plage en médiane, avec une erreur d'ajustement de \num{1.12}\pb, et calibre environ
640 fois plus vite que Fourier en Python. Sur les quatre jours de SPX, il donne les mêmes paramètres et la même
erreur que Fourier (environ \num{29}\pb), en un dixième de seconde au lieu d'une trentaine de secondes. Le coût de la
génération des données et de l'entraînement est amorti dès 5 calibrations face à Fourier en Python, et vers
\num{90} face à Julia : bien moins que ce qu'une banque calibre en une journée.

\paragraph{Limites.} L'erreur sur le SPX vient du modèle, pas du réseau : Heston ne reproduit pas le skew à un mois,
$\xi$ atteint sa borne, et nous avons dû calibrer sur une zone réduite de la surface. Nos données de marché couvrent
quatre jours de novembre 2026, dans un seul régime de volatilité. Enfin, le réseau ne connaît que le domaine sur
lequel il a appris : une surface qui en sort devrait être recalibrée avec Fourier, ou le réseau réentraîné.

\paragraph{Pistes.} Le gain serait bien plus grand pour un modèle sans pricer rapide. Pour les modèles à volatilité
rugueuse \citep{gatheral2018}, qui reproduisent mieux le skew court du SPX, chaque surface demande une simulation de
Monte-Carlo : c'est là que la méthode prend tout son intérêt, comme le montrent \citet{horvath2021}. Une base de
données historique permettrait aussi de tester la calibration sur plusieurs régimes de marché, calmes comme en
crise.
